\documentclass[12pt, letterpaper]{article}

\usepackage[utf8]{inputenc}
\usepackage[T1]{fontenc}
\usepackage[spanish]{babel}
\usepackage[top=2.5cm, bottom=2.5cm, left=3cm, right=3cm]{geometry}
\usepackage{graphicx}
\usepackage{lmodern}
\usepackage{microtype}
\usepackage{hyperref}
\usepackage{listings}
\usepackage{xcolor}
\usepackage{fancyhdr}
\usepackage{titlesec}
\usepackage{enumitem}
\usepackage{booktabs}

% ── COLORES ──────────────────────────────────────────────
\definecolor{primary}{RGB}{0, 70, 127}
\definecolor{codebg}{RGB}{245, 245, 245}
\definecolor{codestr}{RGB}{163, 21, 21}

% ── CÓDIGO ───────────────────────────────────────────────
\lstset{
    backgroundcolor=\color{codebg},
    basicstyle=\ttfamily\small,
    stringstyle=\color{codestr},
    keywordstyle=\color{primary}\bfseries,
    commentstyle=\color{gray}\itshape,
    breaklines=true,
    frame=single,
    rulecolor=\color{gray!40},
    numbers=left,
    numberstyle=\tiny\color{gray},
    xleftmargin=1.5em,
}

% ── SECCIONES ────────────────────────────────────────────
\titleformat{\section}
    {\large\bfseries\color{primary}}
    {\thesection.}{0.5em}{}[\titlerule]

\titleformat{\subsection}
    {\normalsize\bfseries\color{primary!80}}
    {\thesubsection.}{0.5em}{}

% ── ENCABEZADO / PIE ─────────────────────────────────────
\pagestyle{fancy}
\fancyhf{}
\rhead{\textcolor{gray}{\small NOMBRE DEL DOCUMENTO}}
\lhead{\textcolor{gray}{\small \leftmark}}
\rfoot{\textcolor{gray}{\small Página \thepage}}
\lfoot{\textcolor{gray}{\small v1.0 --- \today}}
\renewcommand{\headrulewidth}{0.4pt}

% ── HIPERVÍNCULOS ────────────────────────────────────────
\hypersetup{
    colorlinks=true,
    linkcolor=primary,
    urlcolor=primary,
    citecolor=primary,
}

% ─────────────────────────────────────────────────────────

\begin{document}

% PORTADA
\begin{titlepage}
    \centering
    \vspace*{2cm}
    \includegraphics[width=4cm]{assets/logo.png}\\[2cm]
    {\Huge\bfseries\color{primary} Título del Documento\\[0.5cm]}
    {\Large Subtítulo o versión}\\[2cm]
    \rule{10cm}{0.4pt}\\[0.5cm]
    \begin{tabular}{rl}
        \textbf{Autor:}    & Nombre Autor \\
        \textbf{Empresa:}  & Nombre Empresa \\
        \textbf{Versión:}  & 1.0 \\
        \textbf{Fecha:}    & \today \\
    \end{tabular}\\[2cm]
    \vfill
    {\small Confidencial --- uso interno}
\end{titlepage}

% TABLA DE CONTENIDOS
\tableofcontents
\newpage

% ─────────────────────────────────────────────────────────
\section{Introducción}

Descripción general del documento. Propósito, alcance y audiencia
objetivo.

\subsection{Propósito}

Texto aquí.

\subsection{Alcance}

Texto aquí.

% ─────────────────────────────────────────────────────────
\subsection{Instalación}

\section{Requisitos}

\begin{itemize}[itemsep=2pt]
    \item Requisito uno
    \item Requisito dos
    \item Requisito tres
\end{itemize}

% ─────────────────────────────────────────────────────────
\section{Arquitectura}

\begin{figure}[h]
    \centering
    \includegraphics[width=0.8\textwidth]{assets/diagrama.png}
    \caption{Diagrama de arquitectura}
    \label{fig:arquitectura}
\end{figure}

% ─────────────────────────────────────────────────────────
\section{Instalación y configuración}

\begin{lstlisting}[language=bash]
# Clonar repositorio
git clone https://github.com/org/repo.git
cd repo

# Instalar dependencias
npm install
\end{lstlisting}

\subsection{Variables de entorno}

\begin{center}
\begin{tabular}{lll}
    \toprule
    \textbf{Variable} & \textbf{Descripción} & \textbf{Ejemplo} \\
    \midrule
    \texttt{API\_KEY}   & Clave de acceso a la API & \texttt{abc123}
\\
    \texttt{DB\_URL}    & URL de la base de datos  &
\texttt{postgres://...} \\
    \texttt{PORT}       & Puerto del servidor      & \texttt{3000} \\
    \bottomrule
\end{tabular}
\end{center}

% ─────────────────────────────────────────────────────────
\section{Uso}

\subsection{Ejemplo básico}

\begin{lstlisting}[language=Python]
import modulo

cliente = modulo.Cliente(api_key="xxx")
resultado = cliente.hacer_algo(param="valor")
print(resultado)
\end{lstlisting}

% ─────────────────────────────────────────────────────────
\section{Referencia de API}

\subsection{\texttt{GET /recurso}}

Descripción del endpoint.

\textbf{Parámetros:}
\begin{itemize}
    \item \texttt{id} (requerido): identificador del recurso
\end{itemize}

\textbf{Respuesta:}
\begin{lstlisting}[language=json]
{
"id": 1,
"nombre": "ejemplo",
"estado": "activo"
}
\end{lstlisting}

% ─────────────────────────────────────────────────────────
\section{Solución de problemas}

\begin{description}[style=nextline]
    \item[\texttt{Error: conexión rechazada}]
        Verificar que el servicio esté corriendo y el puerto sea
accesible.
    \item[\texttt{Error: credenciales inválidas}]
        Revisar variables de entorno \texttt{API\_KEY} y
\texttt{SECRET}.
\end{description}

% ─────────────────────────────────────────────────────────
\section{Historial de cambios}

\begin{center}
\begin{tabular}{lll}
    \toprule
    \textbf{Versión} & \textbf{Fecha} & \textbf{Cambios} \\
    \midrule
    1.0 & \today & Versión inicial \\
    \bottomrule
\end{tabular}
\end{center}

\end{document}